\section{Discussion and Conclusion}

\begin{itemize}
    \item The model audit separates basic pattern generation from positional memory. Suitable two-variable Gierer--Meinhardt-type systems already provide a stable resting state, threshold activation, and localized activator peaks.
    \item Source density should therefore not be treated as a generic ingredient that explains regeneration by itself. Its main role is to encode positional information, especially in grafting and polarization experiments.
    \item Prescribed source density acts as an external positional cue, while dynamic source density can provide slow memory. However, dynamic source density alone may produce metastable anchoring rather than genuinely stable heads at arbitrary graft-induced positions.
    \item Introducing a head identity variable with hysteresis separates transient activator peaks from stable morphological commitment. In this interpretation, activator signals are fast, head identity records downstream commitment, and source density is organized by established heads.
    \item Future work should complete the stability and bifurcation analysis, add numerical grafting experiments, calibrate the model biologically, and compare the proposed mechanism with experimental data.
\end{itemize}
